% Hand-written: build.py pastes this into presenting_results.tex, before the results.
\section{Experimental setup}
In this section, we define the experiments that support our theory. Our experiments can be
broadly divided into three categories.

\paragraph{Category 1.}
In the first category, we aim to show that, for any point $x_T$ in the fully noised initial
space, we can predict whether it will lead to a hallucination under the learned score and, if it
does not, which mode it will converge to, without running the full reverse diffusion process.

For our experiments, we consider a Gaussian mixture dataset sampled from $K$ Gaussians:
\begin{equation}
    p_0(x) = \sum_{i=1}^{K} \pi_i \, \mathcal{N}\!\left(x;\, \mu_i,\, \sigma^2 I\right), \qquad \sum_{i=1}^{K} \pi_i = 1, \quad \pi_i > 0,
\end{equation}
where $\mu_i \in \mathbb{R}^d$ denotes the mean of the $i$-th mode, $\pi_i$ its mixing weight,
and $\sigma^2 I$ the (shared) isotropic covariance.

\paragraph{Learned score.}
Let $\mathcal{D}$ denote our sampled data. In phase 1, we train a neural network to learn the
score function of the diffusion process: we generate noisy samples from $\mathcal{D}$ and train
a residual MLP to predict the added noise (DDIM) or the velocity (flow matching). The network
takes $(x, t)$ as input: $x$ is projected linearly to the hidden width and added to a sinusoidal
embedding of $t$ (dimension 128, followed by a linear layer and a LeakyReLU); this passes through
4 residual blocks, each $h \mapsto h + W_2\,\phi(W_1\,\phi(\mathrm{LN}(h)))$ with LayerNorm
$\mathrm{LN}$ and $\phi = \mathrm{LeakyReLU}(0.2)$, and a final LayerNorm, LeakyReLU and linear
layer back to $\mathbb{R}^d$.

The width of the network depends on the dimension of the Gaussian mixture. Concretely, for
$d \le 64$ we use a hidden width of 256, and for $d \in \{128, 256, 512\}$ a width of 1024,
since a fixed width cannot resolve the modes as $d$ grows. We train with Adam (learning rate
$10^{-3}$, cosine-decayed to $10^{-5}$), batch size 512, gradient clipping at 1 and an
exponential moving average of the weights (decay 0.999), on a fixed training set of $10^5$
samples, for $30{,}000\,(1 + d/16)(1 + K/16)$ steps. If the hallucination rate of the trained
sampler exceeds $1.5\%$, the model is retrained once with $1.7\times$ as many steps.

We now describe the classification algorithms.

\paragraph{1) kNN.}
Suppose we have anchor points in the $x_T$ space and $K$ modes. Each anchor point carries a
label in $\{-1\} \cup \{1, \dots, K\}$, where $-1$ is the hallucination class and $k$ is the
$k$-th mode class.

Given a point sampled in the noise space $\mathcal{N}(0, I)$, we first find its $k$ nearest
anchor points in Euclidean distance. For each class, we count how many of these $k$ neighbours
belong to it and take the log of each count (adding a small constant, since $\log 0$ is
undefined). The class with the most votes is the predicted class of the new point.

\paragraph{2) Altered kNN.}
Let $R_{99}$ be the radius of the sphere that contains 99\% of the points sampled from the
Gaussian $\mathcal{N}(\mu_i, \sigma^2 I)$ of mode $i$. The altered kNN then proceeds as follows:

\begin{enumerate}
    \item For each mode $i$ in our dataset, we consider rings at distances
    $r \in \{0.3,\, 0.6,\, 0.9,\, 1.2,\, 1.5\} \times R_{99}$ from the mode centre $\mu_i$.

    \item We assign each ring a weight that decreases linearly with its distance from the mode
    centre ($w \in \{0.84,\, 0.68,\, 0.52,\, 0.36,\, 0.20\}$), and sample $n$ points uniformly
    at random on each ring.

    \item These rings are pulled back with the true score from $x_0$ (the data space) to $x_T$
    (the fully noised space). The pulled-back points are our weighted anchor points.

    \item For any point $x$ in the $x_T$ space, we take its 10 nearest anchor points, denoted
    $\mathcal{K}(x)$, and construct a vector $s(x)$ whose entry for each class $c$ is the
    weighted count of the anchor points of that class:
    \begin{equation}
        s_c(x) = \sum_{j \in \mathcal{K}(x)} w_j \, \mathbf{1}[y_j = c],
    \end{equation}
    where $w_j$ and $y_j$ are the weight and label of anchor point $j$.

    \item We normalise this vector by the total weight of the 10 anchor points,
    \begin{equation}
        \tilde{s}_c(x) = \frac{s_c(x)}{\sum_{j \in \mathcal{K}(x)} w_j},
    \end{equation}
    then divide by a temperature $\tau = 0.1$ and apply a softmax:
    \begin{equation}
        p_c(x) = \frac{\exp\!\left(\tilde{s}_c(x) / \tau\right)}{\sum_{c'} \exp\!\left(\tilde{s}_{c'}(x) / \tau\right)}.
    \end{equation}

    \item Finally, we compute the entropy of $p(x)$ and divide it by $\log K$, the largest
    possible entropy over the $K$ mode classes:
    \begin{equation}
        \hat{H}(x) = -\frac{1}{\log K} \sum_{c} p_c(x) \log p_c(x) \in [0, 1].
    \end{equation}
    This normalised score is 0 for a fully confident prediction and 1 for a uniform one.

    \item The point is predicted to be a hallucination if its confidence $1 - \hat{H}(x)$ falls
    below a threshold, chosen to maximise the hallucination F1 score on 5{,}000 calibration
    points labelled with the true score; otherwise it is assigned to $\arg\max_c p_c(x)$.
\end{enumerate}

\paragraph{3) Quadratic network.}
% TODO

\paragraph{4) Polar network.}
% TODO

\paragraph{Category 2.}
% TODO

\paragraph{Category 3.}
In this category, we present baselines taken from other papers and compare them with our method
on mode prediction accuracy and hallucination F1 score.

\clearpage
